\documentclass[../../main.tex]{subfiles}

\begin{document}

\chapter{Appendix}

Here are the solutions to the practice problems at the end of each
note.

\section{Solutions for Note 1}
\begin{enumerate}
\item
It suffices to show that $\alpha (x, y, z) + \beta (x', y', z') \in S$
for any scalars $\alpha$ and $\beta$. With standard vector addition
and scalar multiplication, $\alpha (x, y, z) + \beta (x', y', z') =
(0, \alpha y + \beta y', \alpha z + \beta z') \in S$. Thus, the set
is a vector space.

\item
Let $f$ and $g$ be functions in $\mathbb{R}$ such that $f'(0) = 0$
and $g'(0) = 0$. Notice that $(f + g)'(0) = f'(0) + g'(0) = 0$ and
$f + g \in S$. Moreover, for any scalars $\alpha$, $(\alpha f)'(0)
= \alpha f'(0) = 0$. Thus $\alpha f \in S$ and $S$ is a vector
space.

\item
First, for any scalars $\alpha$, notice that $\alpha(x_1, \ldots, x_n)
= (\alpha x_1, \ldots, \alpha x_n)$. If each coefficient is divided
by $\alpha \neq 0$, the scalar multiplication holds. The case where
$\alpha = 0$ naturally holds. Continuing, notice that $(x_1, \ldots,
x_n) + (x_1', \ldots, x_n') = (x_1 + x_1', \ldots, x_n + x_n') \in S$
due to the property of addition. Thus, $S$ is a vector space.

\item
Let $V = W_1 \cap \cdots \cap W_n$ where $W_i$ for integer $i \in
[1, n]$ is a subspace of $V$. For $u, v \in V$, notice that since
$u, v \in W_i$ for all $i$ and $\alpha u + \beta v \in W_i$ for any
scalars $\alpha$ and $\beta$ by construction, $\alpha u + \beta v \in
V$ and the intersection of subspaces of a vector space is also a
subspace of the vector space.

\item
First and foremost, it is evident that $e^{a_1 x}$ is linearly
independent since $e^{a_1 x} > 0$ for all $a_1$ and $x$. Therefore,
it suffices to show that if the set is linearly independent for
$n - 1$, then it is linearly independent for $n$. Consider the
following equations for some $c_k$.
\begin{align*}
    \sum_{k=1}^n c_k e^{a_k x} &= 0 \\
    \sum_{k=1}^n c_k a_k e^{a_k x} &= 0 \\
    \sum_{k=1}^n c_k a_n e^{a_k x} &= 0
\end{align*}
Subtracting the third equation, which is obtained by multiplying
$a_n$ in both sides, to the second equation obtained by
differentiating both sides,
\[
\sum_{k=1}^n c_k (a_k - a_n) e^{a_k x}
= \sum_{k=1}^{n-1} c_k (a_k - a_n) e^{a_k x}
= 0
\]
is obtained. Because $a_k \neq a_n$ for integer $k \in [1, n-1]$ and
$\{ e^{a_1 x}, \ldots, e^{a_{n-1} x} \}$ is linearly independent
$c_1 = \cdots = c_{n-1} = 0$. In other words, $c_n e^{a_n x} = 0$ and
$c_n = 0$. Therefore,  $\{ e^{a_1 x}, \ldots, e^{a_n x} \}$ is
linearly independent by induction.

\item
By definition, note that $C = A + B$. Therefore, it suffices to show
that $A \cap B = \{ 0 \}$. Notice that $A$ and $B$ can be written as
$A = \Span\{ (1, 1, 1) \}$ and $B = \Span\{ (20, 15, 12) \}$.
Because $\lambda (1, 1, 1) = (20, 15, 12)$ if and only if $\lambda =
0$, $A \cap B = \{ 0 \}$ and $C = A \oplus B$.

\item
Consider the following equations for some scalars $k_1, k_2, k_3$.
\begin{align*}
    k_1 (a + b) + k_2 (b + c) + k_3 (c + a) &= 0 \\
    a (k_3 + k_1) + b (k_1 + k_2) + c (k_2 + k_3) &= 0
\end{align*}
The linear combinations of $\{ a + b, b + c, c + a \}$ equal zero
if and only if the linear combinations of $a$, $b$, and $c$ equal
zero. Moreover, from the equation above, $k_1 = -k_2 = k_3 = -k_1$.
In other words, $k_1 = k_2 = k_3 = 0$ and $\{ a + b, b + c, c + a \}$
is linearly independent.

\item
Let $f$ and $g$ be polynomials such that $f(0) = g(0) = 0$. Therefore,
$(f + g)(0) = 0$ and $f + g$ is in the set. Moreover, $(\alpha f)(0)
= 0$ for any scalars $\alpha$. Therefore, the set of all polynomials
in $\mathbb{P}^3$ that includes the origin is a subspace of
$\mathbb{P}^3$.
\end{enumerate}
\end{document}


